\chapter*{Anexo 1: Declaración de uso responsable de herramientas de IA}
\addcontentsline{toc}{chapter}{Anexo 1: Declaración de uso responsable de herramientas de IA}

Durante la elaboración de este Trabajo de Fin de Grado se han utilizado herramientas de inteligencia artificial generativa, como modelos de lenguaje de gran escala
y asistentes de programación, de forma auxiliar y bajo criterios de uso ético, crítico y responsable.

El uso de estas herramientas ha estado orientado a apoyar determinadas tareas del proceso de desarrollo del proyecto, sin sustituir la auditoría personal, el criterio
técnico en la toma de decisiones y la responsabilidad acádemica del estudiante. 

En particular, las herramientas de IA se han utilizado para las siguientes finalidades:

\begin{itemize}
    \item Revisión lingüística, mejora de estilo y reformulación de fragmentos redactados previamente por el estudiante.
    \item Generación de ejemplos, explicaciones alternativas o aclaraciones conceptuales.
    \item Apoyo en la depuración, optimización y refactorización del código desarrollado.
    \item Generación de fragmentos de código auxiliares posteriormente revisados, adaptados y validados de acuerdo al criterio del estudiante.
\end{itemize}

En todos los casos expuestos anteriormente, los resultados proporcionados han sido revisados críticamente. El estudiante ha comprobado la corrección técnica, la coherencia
con los objetivos del trabajo, la adecuación al contexto académico y, cuando ha sido necesario, la validez de los algoritmos y fragmentos de código generados.

El uso de IA generativa no ha sustituido el trabajo intelectual propio ni la toma de decisiones fundamentales del proyecto. El diseño, desarrollo, análisis, evaluación, 
redacción final y conclusiones del TFG son responsabilidad del autor de este trabajo.

El estudiante declara que ha utilizado estas herramientas de forma honesta y transparente, evitando presentar como propio trabajo no comprendido, no revisado o generado
automáticamente sin supervisión crítica.

Teniendo esto en cuenta, las herramientas utilizadas han sido: 

\begin{itemize}
    \item \textbf{ChatGPT y Gemini:} utilizados para mejorar la redacción de distintas partes de la memoria, buscando una mayor coherencia y estilo língüístico. También han sido
    utilizados para realizar consultas a nivel técnico sobre cómo utilizar ciertas herramientas (como Locust o Docker), o cómo realizar la conexión entre los componentes.
    
    \item \textbf{Github Copilot:} ha sido utilizado para la generación de fragmentos de código del frontend principalmente, buscando dar una mejora de la usabilidad de la aplicación
    así como del apartado visual. También, ha sido utilizado para optimizar partes del código del despliegue con Docker.
\end{itemize}